\chapter*{Conclusão}
\addcontentsline{toc}{chapter}{Conclusão}
\markboth{CONCLUSÃO}{CONCLUSÃO}

\section{A pergunta e o percurso}

Perguntei como quatro jornais debateram a Caixa de Conversão entre sua criação em 1906 e a suspensão do troco em 1914. Decompus a pergunta para evitar que posição editorial se tornasse um rótulo separado da fonte. Acompanhei objetos concretos, taxa, emissão, lastro, fundos, competências, custódia, crédito e suspensão, e distingui voz do jornal, argumento de terceiro e ato de publicação.

O percurso histórico começou antes da Caixa. A crise da primeira década republicana e o \textit{funding loan} de 1898 deram força à disciplina fiscal, à contração monetária e à valorização. A apreciação do mil-réis, porém, atingiu exportadores numa conjuntura de excesso de café. O Convênio de Taubaté procurou ligar defesa do produto e estabilização cambial. Depois da separação dos projetos, a lei de 1906 produziu um compromisso a 15 pence, reformado para 16 em 1910 e confrontado com a reversão externa a partir de 1912.

Ao longo desse percurso, o debate mudou de composição. Campista e os adversários parlamentares da criação dominaram 1906. Bulhões, Calógeras, Cincinato Braga, Galeão Carvalhal e João Luiz Alves ganharam centralidade na reforma de 1910. Em 1911, a transferência do ouro aos Rothschild deslocou a discussão para custódia e autoridade. Em 1914, liquidez, honra contratual, preservação do lastro e futuro do regime se tornaram inseparáveis da emergência.

\section{O que o material já permite sustentar}

O primeiro resultado é negativo no melhor sentido: a controvérsia não se reduz a metalismo contra papelismo. Em 1906, defensores de 27, de 15 e de 12 pence podiam invocar circulação metálica, conversibilidade e saneamento. Os grupos divergiam sobre paridade, ritmo, alcance e instrumentos. O projeto Campista podia ser chamado de inflacionário por um defensor do par legal e de insuficientemente conversível por um proponente de novo padrão a 12.

O segundo resultado é temporal. As categorias mudaram entre criação, operação, reforma e crise. O \textit{Correio da Manhã} se opôs duramente à taxa de 15 em 1906, registrou a perda de força de 27 em 1907 e defendeu 15 contra 16 em 1910. Uma escala fixa pode descrever isso como travessia de polos, mas a interpretação precisa perguntar se houve revisão doutrinária, adaptação à instituição, defesa de interesses ou continuidade na oposição política.

O terceiro resultado é documental. Publicar uma posição não equivale a adotá-la. \textit{O Paiz} hospedou, em poucos dias de dezembro de 1910, discursos contrários e parecer favorável à elevação. Parte do material estava em seção livre. A multiplicidade é um dado sobre a função do periódico e a composição de suas páginas, mas a posição editorial depende de comentário próprio, enquadramento ou apropriação demonstrável.

O quarto resultado diz respeito a 1914. A crise separou a desejabilidade da estabilidade, a viabilidade da conversibilidade imediata e o custo de defender uma paridade específica. Martim Francisco e Serzedello Corrêa votaram contra a prorrogação da suspensão por razões opostas. A direção do voto não identifica o regime desejado. O primeiro invocava obrigação contratual; o segundo queria permitir que a Caixa esgotasse o ouro e terminasse.

O quinto resultado é a dependência conjuntural. A estabilidade inicial combinou lei, intervenção do Banco do Brasil, financiamento da valorização, exportações e capitais. A reforma de 1910 mostrou que a regra exigia decisão quando o teto foi atingido. A crise de 1913 mostrou o movimento inverso, queda da borracha, problema dos estoques de café, restrição de capitais, perda de reservas e contração. A guerra agravou uma reversão já em curso.

\interpretacaoprovisoria{Esses resultados ainda têm estatuto desigual. A cronologia institucional e os mecanismos apoiados na legislação e na bibliografia são firmes. As sínteses por jornal dependem de completar a leitura comparada, conferir as passagens na imagem e ampliar especialmente 1913 e 1914.}

\section{A posição de cada periódico, matriz para a redação final}

As subseções seguintes não apresentam uma tabela fechada de resultados. Organizam evidências já localizadas e perguntas que precisam ser respondidas antes da versão final.

\subsection{\textit{Correio da Manhã}}

Em 1906, o jornal forneceu a oposição editorial mais explícita entre as peças já lidas. Elogiou a recusa de Rodrigues Alves, apoiou adversários parlamentares de correntes diferentes e atacou Campista e o relatório de Custódio Coelho. Em 1907, registrou o enfraquecimento do horizonte de 27 e considerou 16 uma taxa possível. Em 1908, denunciou a diferença entre a taxa corrente e a usada pela Alfândega. Em 1910, defendeu manter 15 contra a valorização proposta por Bulhões.

As perguntas finais são: que interesses o jornal invocou em cada mudança; até onde a oposição ao governo explica a trajetória; como separou padrão legal, estabilidade e taxa; e o que defendeu durante a saída de ouro e a suspensão. A leitura deve verificar se há continuidade em comércio, lavoura, legalidade ou crítica ao Executivo.

\subsection{\textit{Correio Paulistano}}

O jornal deu grande espaço à justificação técnica do projeto e assumiu em coluna própria a Caixa como parte da política republicana paulista. Sua defesa ligou estabilidade, produção nacional, entrada de capitais e proteção contra especulação. Em 1910, argumentou que o crescimento do ouro e das notas conversíveis aproximava a circulação da metalicidade, exemplo central da combinação entre expansão do meio circulante e finalidade metalista.

As perguntas finais são: quanto da linha resulta do vínculo com o Partido Republicano Paulista; como o jornal tratou os custos e as concessões aos credores da valorização; se sua defesa mudou quando Campista deixou o ministério; e como enquadrou a contração e a suspensão. Também será necessário distinguir os documentos oficiais publicados da voz editorial que os apresentou.

\subsection{\textit{O Paiz}}

O periódico aparece como grande hospedeiro de documentos, discursos e seções livres. Essa função torna seu conteúdo substantivo extenso, mas aumenta o risco de atribuir ao jornal vozes alheias. Um editorial de maio de 1910 oferece voz própria mais segura ao elogiar Murtinho e criticar a assimetria da Caixa. Uma peça de dezembro de 1908 ataca a instituição dentro da luta sucessória, o que recomenda separar hostilidade a Campista de doutrina monetária.

As perguntas finais são: qual foi a linha editorial em 1906; como a proximidade com governos sucessivos alterou o tratamento de Campista, Bulhões e dos presidentes; por que o jornal concedeu espaço a polos opostos em 1910; e que horizonte propôs diante da guerra. A função de hospedagem precisa ser descrita como resultado próprio, não tratada apenas como ruído.

\subsection{\textit{Gazeta de Notícias}}

Em março de 1906, o jornal apoiou a estabilização, criticou o desenho de duas circulações e procurou combinações entre 12 e 27. Em novembro, declarou apoio decidido ao projeto aprovado. Em janeiro de 1912, afirmou retrospectivamente sua responsabilidade no auxílio à Caixa e procurou normalizar retiradas de ouro. A falta de objetos de 1913 impede acompanhar o momento da reversão nas mesmas condições dos outros títulos. Em 1914, o jornal reproduziu o debate parlamentar, inclusive votos iguais por razões incompatíveis.

As perguntas finais são: como ocorreu a passagem entre crítica do desenho e apoio à lei; qual foi sua posição na elevação para 16; como sua autodescrição de 1912 corresponde ao arquivo anterior; e se, em 1914, houve apropriação editorial das falas reproduzidas. A lacuna de 1913 deve permanecer visível na comparação.

\section{Credores e banqueiros estrangeiros}

O material permite formular uma questão transversal. Os Rothschild e outros banqueiros aparecem como agentes do governo, credores, destinatários de telegramas e depositários de ouro. São designados por cargo ou perífrase em 456 das 8.331 páginas que mencionam a Caixa, contra 170 que trazem o nome de um Rothschild, e 418 páginas designam o credor sem nunca nomeá-lo, enquanto o parlamento quase não o nomeia. O contraste sugere que a imprensa acompanhava a operação cotidiana do crédito internacional com uma linguagem que naturalizava a presença do agente financeiro.

A interpretação final deverá evitar dois extremos. O primeiro é tornar banqueiros estrangeiros causa única das decisões brasileiras. O segundo é tratá-los como intermediários neutros. A negociação de 1898, a garantia da valorização, a custódia de 1911, as dificuldades de 1913 e o novo \textit{funding} de 1914 mostram uma relação assimétrica na qual governo, estados, exportadores e credores tinham interesses próprios e capacidades diferentes.

\section{Contribuição, limites e agenda imediata}

A contribuição principal é reconstruir o debate em sua diversidade documental. O censo digital oferece escala, mas a distinção entre rotina e argumento reduz drasticamente o material interpretável. A descrição inventarial de gêneros, vozes, atores, termos e episódios torna possível escrever sem depender de uma pontuação única de posicionamento. Também preserva o caminho entre uma afirmação e a página que a sustenta.

O primeiro limite é o acervo, especialmente a falta da \textit{Gazeta de Notícias} em 1913 e a ausência serial do diário do \textit{Jornal do Commercio}. O segundo é técnico, pois o ruído de OCR varia entre jornais e anos, e a segmentação automática pode misturar colunas. O terceiro é interpretativo, já que poucas peças editoriais foram identificadas em algumas fases. O quarto é o estado do trabalho, com citações primárias ainda pendentes de conferência visual.

\interpretacaoprovisoria{A resposta mais promissora à pergunta da dissertação é que os jornais não apenas ocuparam posições num debate estável. Eles ajudaram a mudar sua composição. Em 1906 organizaram alternativas de paridade e desenho; nos primeiros anos avaliaram efeitos e incoerências; em 1910 reabriram taxa e autoridade; na crise separaram contrato, lastro, liquidez e futuro do regime. Comparar jornais significa acompanhar essas transformações, não calcular uma média atemporal.}

A agenda imediata está concentrada. É preciso conferir as citações na imagem, concluir as janelas de 1906, 1910 e 1914, reconstruir a cronologia mensal de 1912 a 1914 e preencher a matriz de cada periódico com evidência de voz. Feito isso, a conclusão poderá substituir perguntas por respostas e distinguir com clareza o que confirma, corrige ou complica a historiografia.
